
\chapter{CONCLUSIONES}

La presente investigación abordó el desarrollo de un algoritmo de clasificación basado en
aprendizaje profundo para la detección de retinopatía diabética en imágenes del fondo del
ojo, partiendo de una revisión exhaustiva de la literatura científica sobre técnicas de
aprendizaje profundo aplicadas a la detección de retinopatía diabética, lo que permitió
identificar las arquitecturas más prometedoras y las estrategias de preprocesamiento más
efectivas para este dominio. Este análisis fundamentó las decisiones metodológicas
adoptadas y contextualizó los resultados obtenidos dentro del estado del arte.

Se recopiló y organizó un conjunto de datos clínicos compuesto por múltiples fuentes: el
dataset público de Kaggle con 35{,}121 imágenes para entrenamiento, complementado con los
conjuntos MESSIDOR-2, IDRiD y FOSCAL para validación externa. La implementación de
técnicas de preprocesamiento ---incluyendo normalización del color mediante filtro
gaussiano, enmascaramiento periférico y balanceo de clases--- demostró ser determinante
para optimizar la calidad de los datos de entrada.

Se desarrolló un algoritmo de clasificación basado en redes neuronales convolucionales
utilizando aprendizaje por transferencia con cuatro arquitecturas preentrenadas:
MobileNetV2, ResNet50V2, InceptionV3 y VGG19. La arquitectura propuesta incorporó bloques
Multi-Head Attention y capas GaussianNoise como estrategias de regularización, logrando un
equilibrio efectivo entre capacidad de generalización y eficiencia computacional.

Se evaluó el desempeño del sistema mediante comparación con trabajos previos del estado
del arte y mediante validación externa, sin reentrenamiento ni ajuste de pesos, en tres
conjuntos de datos independientes al de entrenamiento. El modelo MobileNetV2 con filtrado alcanzó
una sensibilidad del 96.5\% y un F1-score de 0.96 en clasificación binaria, con intervalos
de confianza del 95\% que confirman la robustez estadística de estos resultados (AUC-ROC:
0.9636--0.9936). En clasificación multiclase, se determinó que el canal verde proporciona
el mejor desempeño (exactitud de 0.55), hallazgo que posee fundamento fisiológico dada la
absorción diferencial de la hemoglobina en esta longitud de onda.

Los resultados demuestran que la implementación de modelos de aprendizaje profundo con
arquitecturas ligeras y preprocesamiento robusto constituye una alternativa técnicamente
viable para el tamizaje automatizado de retinopatía diabética. Esta solución posee
potencial para abordar la escasez de especialistas en oftalmología documentada en
Colombia, donde la proporción actual de 1.56 oftalmólogos por cada 100{,}000 habitantes
resulta insuficiente para atender a los aproximadamente 1.2 millones de personas con
retinopatía diabética en el país.

Se concluye que los algoritmos de aprendizaje profundo, particularmente cuando se combinan
con técnicas de aprendizaje por transferencia y preprocesamiento especializado, ofrecen un
enfoque prometedor para la detección temprana de retinopatía diabética, con métricas de
desempeño comparables a las reportadas en la literatura internacional. La evaluación en
conjuntos de datos externos, realizada sin reentrenamiento ni ajuste de pesos, mostró
resultados alentadores pero heterogéneos entre las distintas fuentes evaluadas, lo que
indica una capacidad de generalización todavía parcial y sensible a las diferencias de
adquisición, población y protocolo de cada institución. En ausencia de una evaluación
prospectiva y de una validación clínica formal, estos hallazgos deben interpretarse como
evidencia preliminar del potencial del sistema como herramienta de apoyo al tamizaje, y no
como una demostración de aplicabilidad clínica ya establecida; su eventual implementación
en un entorno real requeriría estudios adicionales de validación prospectiva y evaluación
conjunta con especialistas en oftalmología.
